\section{Guia de capturas}

Las capturas son opcionales para compilar el manual, pero recomendadas para
compartirlo con alumnos. Guardarlas dentro de \texttt{figures/} usando los
nombres exactos.

\begin{longtable}{p{0.28\linewidth}p{0.64\linewidth}}
\toprule
Archivo & Momento de captura \\
\midrule
\texttt{01-docker-servicios.png} & Terminal con contenedores activos. \\
\texttt{02-health-ready.png} & Health check del backend. \\
\texttt{03-pacientes-listado.png} & Pantalla \texttt{/patients}. \\
\texttt{04-sandbox-sesion.png} & Menu superior del sandbox. \\
\texttt{05-ficha-paciente.png} & Ficha de un paciente. \\
\texttt{06-editar-datos.png} & Drawer \texttt{Editar datos del paciente}. \\
\texttt{07-reglas-catalogo.png} & Catalogo de reglas CQL. \\
\texttt{08-editor-cql.png} & Editor CQL integrado. \\
\texttt{09-validacion-elm.png} & Resultado de validacion y ELM. \\
\texttt{10-prueba-no-aplica.png} & Prueba de regla que no aplica. \\
\texttt{11-card-activa.png} & Card CDS visible en ficha. \\
\texttt{12-actividad-cds.png} & Pantalla \texttt{/activity}. \\
\texttt{13-trazabilidad-hook.png} & Detalle visual de la ejecucion CDS Hooks. \\
\texttt{14-cds-discovery.png} & Discovery CDS Hooks por terminal. \\
\texttt{15-cds-post-card.png} & POST CDS Hooks con respuesta de cards. \\
\bottomrule
\end{longtable}

Para una guia aun mas clara, tambien se pueden agregar:

\begin{itemize}
  \item \texttt{16-error-cql.png}: error de CQL en el editor.
  \item \texttt{17-dos-navegadores.png}: dos sandboxes distintos.
  \item \texttt{18-reset-sandbox.png}: reinicio de sandbox.
\end{itemize}
